\chapter{Deferred Proofs}\label{app:proofs}

This appendix collects those proofs which are given in full but whose length would have interrupted
the argument in the main text. The organising principle is the complement of the one governing
Appendix \ref{app:auxiliary}: a result whose proof is short enough to be read in passing is proved
where it is stated, a result we merely invoke is recorded without proof in Appendix
\ref{app:auxiliary}, and a result we genuinely use but whose proof runs to more than a few lines is
stated in the body and proved here.

Each section below carries a pointer back to the statement it establishes, and a reference to
\cite{Kuo2006}, whose treatment the arguments follow.

\section{Doob's submartingale inequality}\label{app:proof:doobmaximal}
We prove Proposition \ref{prop:doobmaximal} of Chapter \ref{ch:cap2}.

\textbf{\textit{Proof:}} Let $Q$ be an enumeration of the rationals in $[0,T]$. Right-continuity gives
$\sup_{0\leq t\leq T} Z_t = \sup_{r \in Q} Z_r$ almost surely, so it suffices to bound the maximum
over a finite set $\{r_1 < \cdots < r_k\}$ and let $k \to \infty$. On such a set, decompose
$\{\max_\nu Z_{r_\nu} \geq \varepsilon\}$ into the disjoint events $A_\nu$ on which $r_\nu$ is the
\emph{first} index at which the level is reached. Each $A_\nu \in \mathcal{F}_{r_\nu}$, and on
$A_\nu$ we have $Z_{r_\nu} \geq \varepsilon$, so the submartingale property gives
\[
\varepsilon\,\mathbb{P}(A_\nu) \;\leq\; \mathbb{E}\bigl[Z_{r_\nu}\mathds{1}_{A_\nu}\bigr]
\;\leq\; \mathbb{E}\bigl[\mathbb{E}[Z_T \mid \mathcal{F}_{r_\nu}]\,\mathds{1}_{A_\nu}\bigr]
\;=\; \mathbb{E}\bigl[Z_T\,\mathds{1}_{A_\nu}\bigr] \;\leq\; \mathbb{E}\bigl[\max\{Z_T,0\}\,\mathds{1}_{A_\nu}\bigr].
\]
Summing over $\nu$ and using disjointness bounds $\varepsilon\,\mathbb{P}(\max_\nu Z_{r_\nu}\geq\varepsilon)$
by $\mathbb{E}[\max\{Z_T,0\}]$, uniformly in $k$. Letting $k \to \infty$ gives the claim. The
martingale case follows by applying this to the submartingale $|X_t|$, which is one by the
conditional Jensen inequality. \hfill $\blacksquare$

\section{Density of step processes in $L^2_{ad}([0,T]\times\Omega)$}\label{app:proof:density}
We prove Lemma \ref{lem:density} of Chapter \ref{ch:cap2}.

\textbf{\textit{Proof:}} The argument proceeds by successive approximation in three steps, each
enlarging the class of integrands. Write $\lVert f\rVert^2 := \int_0^T\mathbb{E}[f(t)^2]\,dt$ for the
norm of $L^2_{ad}([0,T]\times\Omega)$, and recall that it suffices to approximate an arbitrary
$f$ within $\varepsilon$ by a step process, since the step processes then form a dense subspace.

\medskip
\noindent\textbf{Step 1: $f$ bounded with continuous sample paths.} Let $|f| \leq M$ and let
$\mathcal{P}_n = \{0 = t^n_0 < \cdots < t^n_{k_n} = T\}$ be partitions with
$\lVert\mathcal{P}_n\rVert \to 0$. Put
\[
f_n(t) := \sum_{i=1}^{k_n} f(t^n_{i-1})\,\mathds{1}_{[t^n_{i-1},\,t^n_i)}(t).
\]
Each $f_n$ is a step process, and it is \emph{adapted} because the coefficient $f(t^n_{i-1})$ is
$\mathcal{F}_{t^n_{i-1}}$-measurable and multiplies the indicator of an interval beginning at
$t^n_{i-1}$ --- this is the left-endpoint choice, and it is the only reason adaptedness survives. For
almost every $\omega$ the path $t \mapsto f(t,\omega)$ is continuous on the compact $[0,T]$, hence
uniformly continuous, so $f_n(t,\omega) \to f(t,\omega)$ for every $t$. Since
$|f_n - f|^2 \leq 4M^2$, which is integrable on $[0,T]\times\Omega$ (a finite measure space),
dominated convergence gives $\lVert f_n - f\rVert \to 0$.

\medskip
\noindent\textbf{Step 2: $f$ bounded, sample paths not assumed continuous.} Let $|f| \leq M$. Extend
$f$ by $0$ outside $[0,T]$ and mollify in time: for $n \in \mathbb{N}$ set
\[
g_n(t) := n\int_{t - 1/n}^{t} f(s)\,ds = \int_0^1 f\bigl(t - \tfrac{u}{n}\bigr)du .
\]
Each $g_n$ has continuous sample paths, since it is an indefinite integral, and $|g_n| \leq M$.
Adaptedness is preserved because $g_n(t)$ averages $f$ over $[t-1/n,\,t]$ only, using no information
after $t$; this is where averaging \emph{backwards} rather than symmetrically matters. By the
Lebesgue differentiation theorem, $g_n(t,\omega) \to f(t,\omega)$ for almost every $t$, for almost
every $\omega$, and dominated convergence again gives $\lVert g_n - f\rVert \to 0$. Step 1 now
approximates each $g_n$ by step processes, and a diagonal argument produces step processes converging
to $f$.

\medskip
\noindent\textbf{Step 3: general $f \in L^2_{ad}([0,T]\times\Omega)$.} Truncate at level $n$ by
setting $f^{(n)} := (f \wedge n) \vee (-n)$. Each $f^{(n)}$ is bounded and remains adapted, being a
composition of $f$ with a continuous deterministic function. Moreover $|f^{(n)} - f|^2 \leq f^2$
pointwise and $f^{(n)} \to f$ pointwise, so dominated convergence --- with dominating function $f^2$,
integrable precisely because $f \in L^2_{ad}$ --- yields $\lVert f^{(n)} - f\rVert \to 0$. Combining
with Step 2 completes the proof.

\medskip
\noindent Adaptedness is the only delicate point throughout, and it is worth noting that each of the
three steps preserves it for a different reason: the left-endpoint evaluation in Step 1, the backward
averaging window in Step 2, and the deterministic truncation in Step 3. \hfill $\blacksquare$

\section{Continuity of the sample paths of an Itô integral}\label{app:proof:itocontinuity}
We prove Theorem \ref{thm:itocontinuity} of Chapter \ref{ch:cap2}.

\textbf{\textit{Proof:}} Let $\{f_n\}$ be step processes converging to $f$ in
$L^2_{ad}([0,T]\times\Omega)$; the associated processes $X^n_t = \int_0^t f_n\,dW$ have continuous
paths by construction, being finite sums of continuous terms. By Theorem
\ref{theorem:martingale_property} the difference $X^n - X^m$ is a martingale, so Doob's $L^2$
maximal inequality (Proposition \ref{prop:dooblp}) together with Itô's isometry gives
\[
\mathbb{E}\Bigl[\sup_{0\leq t\leq T}\bigl|X^n_t - X^m_t\bigr|^2\Bigr]
\;\leq\; 4\,\mathbb{E}\bigl[|X^n_T - X^m_T|^2\bigr]
\;=\; 4\int_0^T \mathbb{E}\bigl[(f_n - f_m)^2\bigr]\,dt \longrightarrow 0 .
\]
Passing to a subsequence for which the right-hand side is summable, the Borel--Cantelli lemma yields
$\sup_{t\leq T}|X^{n_k}_t - X_t| \to 0$ almost surely. A uniform limit of continuous functions is
continuous, so almost all paths of $X$ are continuous on $[0,T]$. \hfill $\blacksquare$

\section{The Itô integral as a limit of Riemann sums}\label{app:proof:riemannlimit}
We prove Theorem \ref{theorem:riemannian_limit} of Chapter \ref{ch:cap2}.

\textbf{\textit{Proof:}} Let $f_\mathcal{P}(t) = \sum_i f(t_{i-1})\mathds{1}_{[t_{i-1},t_i)}(t)$ be
the left-endpoint step approximation, so that the Riemann sum is exactly $I(f_\mathcal{P})$. By Itô's
isometry it suffices to show that $\int_0^T \mathbb{E}[(f - f_\mathcal{P})^2]\,dt \to 0$. Expanding
the square, this integral is a sum of terms of the form
$\mathbb{E}[f(t)^2] - 2\mathbb{E}[f(t)f(t_{i-1})] + \mathbb{E}[f(t_{i-1})^2]$, each built from the
function $(s,t)\mapsto\mathbb{E}[f(t)f(s)]$ evaluated at points within $\lVert\mathcal{P}\rVert$ of
the diagonal. That function is continuous by hypothesis, hence uniformly continuous on the compact
$[0,T]^2$, so the whole expression tends to $0$ uniformly as $\lVert\mathcal{P}\rVert \to 0$. \hfill $\blacksquare$

\section{Lévy's characterisation theorem}\label{app:proof:levy}
We prove Theorem \ref{thm:levycharacterisation} of Chapter \ref{ch:cap2}.

\textbf{\textit{Proof:}} The one-dimensional case is proved by applying Itô's rule to
$t \mapsto e^{iuX_t}$ for fixed $u \in \mathbb{R}$. Since $X$ is a continuous local martingale with
$\langle X\rangle_t = t$, the second-order term contributes $-\tfrac12 u^2 e^{iuX_t}\,dt$, and after
localisation and taking conditional expectations one obtains, for $s<t$,
\[
\mathbb{E}\bigl[e^{iu(X_t - X_s)} \mid \mathcal{F}_s\bigr]
= 1 - \tfrac12 u^2 \int_s^t \mathbb{E}\bigl[e^{iu(X_r-X_s)}\mid\mathcal{F}_s\bigr]\,dr .
\]
This is \emph{not} yet a deterministic equation: for each $r$ the quantity
$\mathbb{E}[e^{iu(X_r-X_s)}\mid\mathcal{F}_s]$ is an $\mathcal{F}_s$-measurable random variable, so
the unknown is a random function and no uniqueness theorem for integral equations applies to it as it
stands. The standard device removes the randomness by testing. Fix $A \in \mathcal{F}_s$, multiply
through by $\mathds{1}_A$, take expectations, and interchange expectation and the $dr$ integral,
which is legitimate by Fubini since the integrand has modulus one. Writing
\[
h_A(t) := \mathbb{E}\bigl[\mathds{1}_A\,e^{iu(X_t-X_s)}\bigr],
\]
we obtain the genuinely deterministic linear integral equation
\[
h_A(t) = \mathbb{P}(A) - \tfrac12 u^2 \int_s^t h_A(r)\,dr, \qquad t \geq s,
\]
whose unique solution is $h_A(t) = \mathbb{P}(A)\,e^{-u^2(t-s)/2}$. Since this holds for every
$A \in \mathcal{F}_s$ and the right-hand side is the integral over $A$ of the constant
$e^{-u^2(t-s)/2}$, we conclude
\[
\mathbb{E}\bigl[e^{iu(X_t-X_s)} \mid \mathcal{F}_s\bigr] = e^{-u^2(t-s)/2} \qquad \text{almost surely.}
\]
The conditional characteristic function of $X_t - X_s$ given $\mathcal{F}_s$ is therefore
deterministic and equal to that of $\mathcal{N}(0,t-s)$, which is simultaneously
independence of the increment from $\mathcal{F}_s$ and the correct Gaussian law; Proposition
\ref{prop:independent-increments} then supplies independence of increments. The multidimensional case
is identical, applied to $e^{iu\transp X_t}$ with the cross-variations $\langle X^i,X^j\rangle_t
= \delta_{ij}t$ making the exponent $-\tfrac12\lVert u\rVert^2(t-s)$. \hfill $\blacksquare$

\section{Girsanov's theorem}\label{app:proof:girsanov}
We prove Theorem \ref{girsanovtheorem} of Chapter \ref{ch:cap2}.

\textbf{\textit{Proof:}} By Theorem \ref{thm:levycharacterisation} it suffices to show that, under
$\mathbb{Q}$, the process $Z$ is a continuous local martingale null at the origin with
$\langle Z\rangle_t = t$. Continuity and $Z_0=0$ are clear.

For the martingale property we use Lemma \ref{lem:qmartingale}, which reduces the claim to showing that
$Z(t)\mathcal{E}_\alpha(t)$ is a $\mathbb{P}$-martingale. Itô's product formula gives
\[
d\bigl(Z_t\mathcal{E}_\alpha(t)\bigr)
= \mathcal{E}_\alpha(t)\,dZ_t + Z_t\,d\mathcal{E}_\alpha(t) + dZ_t\,d\mathcal{E}_\alpha(t) .
\]
Now $dZ_t = dW_t - \alpha(t)\,dt$ and $d\mathcal{E}_\alpha(t) = \alpha(t)\mathcal{E}_\alpha(t)\,dW_t$,
so by the Itô table $dZ_t\,d\mathcal{E}_\alpha(t) = \alpha(t)\mathcal{E}_\alpha(t)\,dt$. Substituting,
the two $dt$ terms cancel exactly:
\[
d\bigl(Z_t\mathcal{E}_\alpha(t)\bigr)
= \mathcal{E}_\alpha(t)\bigl(dW_t - \alpha(t)dt\bigr) + Z_t\alpha(t)\mathcal{E}_\alpha(t)\,dW_t
+ \alpha(t)\mathcal{E}_\alpha(t)\,dt
= \bigl(1 + Z_t\alpha(t)\bigr)\mathcal{E}_\alpha(t)\,dW_t .
\]
The right-hand side has no drift, so $Z\mathcal{E}_\alpha$ is a $\mathbb{P}$-local martingale, and
under the stated integrability it is a martingale; hence $Z$ is a $\mathbb{Q}$-martingale.

Finally, $\langle Z\rangle_t = \langle W\rangle_t = t$, because subtracting the finite-variation term
$\int_0^t\alpha(s)ds$ does not change the quadratic variation, and because quadratic variation is a
pathwise limit and is therefore unaffected by the passage to the equivalent measure $\mathbb{Q}$. The
same argument applied to $Z^2_t - t$ confirms the compensator. Lévy's theorem now identifies $Z$ as a
$\mathbb{Q}$-Brownian motion. \hfill $\blacksquare$
